\documentclass[11pt]{article}
\usepackage{../../../tex/lecturestyle}

\title{Parking Function}
\author{}
\date{}

\begin{document}
\thispagestyle{fancy}

\lecturetitleband{Combinatorics}{Parking Function}{4}

\section{Overview}
A \emph{parking function} is a combinatorial object with a memorable
story: $n$ cars, each with a preferred parking spot among $n$ numbered
spots, try to park one at a time, each driver taking the first free spot
at or after their preference (wrapping is not allowed --- if no such spot
exists, parking fails outright). The sequence of preferences is a
\emph{parking function} exactly when every car succeeds. The count of
such sequences has a strikingly clean closed form, $(n+1)^{n-1}$, proved
here by a clever trick: extend the parking lot into a circle.

\section{Definition}
\begin{definitionbox}
Fix $n \ge 1$ and a sequence $a = (a_1, \ldots, a_n)$ with each
$a_i \in \{1, \ldots, n\}$. Process cars $i = 1, \ldots, n$ in order: car
$i$ scans spots $a_i, a_i+1, \ldots, n$ and parks in the first one still
free, or the whole process fails if none is free. $a$ is a \emph{parking
function} if every car parks successfully.
\end{definitionbox}

\begin{examplebox}
For $n=5$, $(3,2,2,4,1)$ is a parking function: cars take spots
$3,2,4,5,1$ respectively (car $3$ wants spot $2$ but finds it taken, so
takes spot $4$; and so on). By contrast $(3,4,3,4,1)$ fails: by the time
the fourth car (wanting spot $4$) arrives, spots $1,3,4$ are already
taken by earlier cars, and no free spot $\ge 4$ remains.
\end{examplebox}

% FIGURE: pf_street
\begin{figure}[H]
\centering
\begin{tikzpicture}[x=1.45cm,y=1.0cm]
\node[lcell,minimum width=1.3cm,minimum height=0.9cm,lhi] (p1) at (1,0) {car 2};
\node[llabel] at (1,-0.75) {spot 1};
\node[lcell,minimum width=1.3cm,minimum height=0.9cm,lhi] (p2) at (2,0) {car 1};
\node[llabel] at (2,-0.75) {spot 2};
\node[lcell,minimum width=1.3cm,minimum height=0.9cm,lhi] (p3) at (3,0) {car 3};
\node[llabel] at (3,-0.75) {spot 3};
\node[lcell,minimum width=1.3cm,minimum height=0.9cm,lhi] (p4) at (4,0) {car 4};
\node[llabel] at (4,-0.75) {spot 4};
\node[lcell,minimum width=1.3cm,minimum height=0.9cm,lhi] (p5) at (5,0) {car 5};
\node[llabel] at (5,-0.75) {spot 5};
\node[llabel,text=lectureaccent] at (2,1.0) {};
\node[llabel,text=lectureaccent] at (1,1.0) {};
\draw[larrow,lecturewarn,bend left=45] ([yshift=0.0cm]p2.north) to node[above,llabel] {car 3} (p3.north);
\draw[larrow,lecturewarn,bend left=45] ([yshift=0.0cm]p1.north) to node[above,llabel] {car 4} (p4.north);
\node[llabel,text=lectureaccent] at (5,1.0) {};
\node[llabel,anchor=west,align=left] at (0.2,-1.55) {preferences $(a_1,\dots,a_5)=(2,1,2,1,5)$: cars 1, 2, 5 park at their preferred spot;\\ car 3 prefers 2 (taken) and moves on to 3; car 4 prefers 1 and drives on to 4.};
\end{tikzpicture}
\caption{Cars park one at a time; each takes its preferred spot or, if that is taken, the first free spot after it. Brown arcs show the cars that had to move past their preference. All five cars park, so $(2,1,2,1,5)$ is a parking function: its sorted entries $(1,1,2,2,5)$ satisfy $b_i\le i$.}\label{fig:pf_street}
\end{figure}

\section{Equivalent Characterizations}
\begin{theorembox}
For $a = (a_1,\ldots,a_n)$ with values in $\{1,\ldots,n\}$, the following
are equivalent:
\begin{enumerate}[leftmargin=1.5em]
  \item $a$ is a parking function.
  \item For every $k=1,\ldots,n$, at least $k$ of the $a_i$ satisfy
        $a_i \le k$.
  \item Sorting $a$ into $b_1 \le \cdots \le b_n$, every $b_k \le k$.
  \item There is a permutation $p$ of $\{1,\ldots,n\}$ with $a_i \le p_i$
        for every $i$ (take $p_i$ to be the spot car $i$ actually parks
        in).
\end{enumerate}

\emph{Proof sketch.} (2)$\Leftrightarrow$(3) is immediate (re-sorting).
(1)$\Rightarrow$(4): let $p_i$ be where car $i$ parks; it never parks
before its preference, so $a_i \le p_i$, and no two cars share a spot, so
$p$ is a permutation. (4)$\Rightarrow$(2): among the $k$ indices $i$ with
$p_i \in \{1,\ldots,k\}$, each has $a_i \le p_i \le k$, giving at least
$k$ such indices. (2)$\Rightarrow$(1) by induction on $n$: condition (2)
guarantees some $a_{i_0}=1$; if any later $a_i=1$ too, changing it to $2$
preserves both parking-function-ness (spot $1$ is already taken by car
$i_0$ when car $i$ arrives, so its behavior doesn't change) and condition
(2) (only the $k=1$ count could change, and it doesn't, by choice of
$i_0$ as the first index with $a_i=1$); so WLOG $i_0$ is the only index
with value $1$. Deleting $a_{i_0}$ and subtracting $1$ from every
remaining entry yields a length-$(n{-}1)$ sequence that is a parking
function (in the smaller sense) exactly when the original is, and
satisfies condition (2) exactly when the original does --- so induction
applies. \qed
\end{theorembox}

\begin{notebox}
Condition (4) is a perfect matching statement, and the whole equivalence
can also be derived from Hall's marriage theorem; the parking procedure
itself is essentially a greedy algorithm for finding such a matching.
\end{notebox}

\section{Counting Parking Functions}
\begin{theorembox}
The number of parking functions of length $n$ is $(n+1)^{n-1}$.

\emph{Proof.} Extend the lot to a \emph{circle} of $n+1$ spots numbered
$0, \ldots, n$, and allow preferences $a_i \in \{0, \ldots, n\}$: car $i$
now scans spots $a_i, a_i{+}1, \ldots, n, 0, 1, \ldots, a_i{-}1$ ---
every spot, circularly --- so with only $n$ cars and $n+1$ spots, parking
\emph{never} fails, and exactly one spot is left empty at the end.

There are $(n+1)^n$ sequences $a \in \{0,\ldots,n\}^n$ in total. By the
rotational symmetry of the circle (relabeling every spot by a fixed
rotation just relabels which sequences leave which spot empty, without
changing how many leave \emph{some} spot empty), the number of sequences
leaving spot $s$ empty is the same for every $s \in \{0,\ldots,n\}$,
hence equals $(n+1)^n / (n+1) = (n+1)^{n-1}$.

It remains to identify ``sequences leaving spot $0$ empty'' with ordinary
parking functions. If some $a_i = 0$, car $i$ finds spot $0$ free at its
turn (nothing fills spot $0$ before any car explicitly aiming for it, as
the circular scan reaches $0$ only after wrapping past $n$) and takes it
immediately, so spot $0$ ends up occupied. So spot $0$ can only remain
empty when every $a_i \in \{1,\ldots,n\}$ --- and for such $a$, the
circular process, restricted to spots $\{1,\ldots,n\}$, behaves exactly
like the ordinary (non-circular) parking process: a car only ever reaches
spot $0$ via wraparound, which happens exactly when it would have failed
to park in $\{1,\ldots,n\}$ under the ordinary rules. So spot $0$ stays
empty exactly when every car parks within $\{1,\ldots,n\}$ under the
ordinary process --- i.e.\ exactly when $a$ is an ordinary parking
function. This identifies the two counts, giving $(n+1)^{n-1}$. \qed
\end{theorembox}

\subsection{A generalization}
The same circular argument extends to sequences shorter than the lot.

\begin{theorembox}
For $1 \le m \le n$, the number of length-$m$ sequences
$a=(a_1,\ldots,a_m)$ with values in $\{1,\ldots,n\}$ such that, for every
$k=1,\ldots,n$, at least $m-n+k$ of the $a_i$ satisfy $a_i \le k$, is
\[
  (n+1-m)\,(n+1)^{m-1}.
\]
(The proof mirrors the one above, with $m$ cars circulating among $n+1$
spots; each of the $n+1-m$ ways to leave $n+1-m$ spots simultaneously
untouched contributes equally by the same symmetry argument.)
\end{theorembox}

\section{Reference Implementations}
Both examples below implement the direct simulation, the threshold
condition, and the circular extension, cross-checking all three against
each other (and against Proposition 5's generalized count) by brute
force, plus reproducing the lecture's own worked examples exactly.

\codenote
\subsection{Rust}
\corelisting{Rust}{example.rs}

\subsection{Python}
\corelisting{Python}{example.py}

\section{Worked Example}
\begin{examplebox}
Among the four sequences $(4,1,5,1,3)$, $(5,3,1,4,3)$, $(4,1,6,5,4,1)$,
$(3,6,1,5,3,1)$, exactly the first and the last are parking functions ---
confirmed both by direct simulation and by the threshold condition in
both reference implementations. For $n=1,2,3,4,5$, brute force confirms
$1, 3, 16, 125, 1296$ parking functions respectively, matching
$(n+1)^{n-1}$.
\end{examplebox}

\section{Further Reading}
\begin{itemize}[leftmargin=1.5em]
  \item AtCoder Algorithm Lectures index (source of this curriculum):
        \url{https://info.atcoder.jp/entry/algorithm_lectures/index}
  \item A. G. Konheim, B. Weiss, ``An Occupancy Discipline and
        Applications'', SIAM Journal on Applied Mathematics 14(6), 1966,
        pp.\ 1266--1274.
  \item R. P. Stanley, \emph{Parking functions}:
        \url{https://math.mit.edu/~rstan/transparencies/parking.pdf}
\end{itemize}

\end{document}
